% =============================================================================
%  CHAPTER 6 -- MEAN REVERSION AND THE ORNSTEIN--UHLENBECK PROCESS
% =============================================================================
\section{Mean reversion and the Ornstein--Uhlenbeck process}
\label{ch:ou}

Chapter~\ref{ch:timeseries} established \emph{whether} a spread is stationary. This
chapter quantifies \emph{how fast} it comes home, which is the difference between a
tradeable signal and an interesting curiosity.

\subsection{Stochastic differential equations in one page}
\label{sec:ou:sde}

A continuous-time diffusion is written
\begin{equation}
  \dd X_t \;=\; \underbrace{a(X_t,t)}_{\text{drift}}\dd t
  \;+\; \underbrace{b(X_t,t)}_{\text{diffusion}}\dd W_t ,
  \label{eq:sde}
\end{equation}
where $W_t$ is a standard Brownian motion (Wiener process): $W_0=0$, independent
increments, and $W_{t+s}-W_t\sim\mathcal N(0,s)$. The notation $\dd W_t$ is formal ---
Brownian paths are nowhere differentiable, so \eqref{eq:sde} means the integral equation
$X_t = X_0+\int_0^t a\dd s + \int_0^t b\dd W_s$ with the second integral in the Itô sense.

Two properties of Brownian motion do all the work:
\begin{align}
  \E\Bigl[\int_0^t f\dd W_s\Bigr] &= 0 &&\text{(the Itô integral is a martingale)}, \label{eq:ito:mean}\\
  \Var\Bigl[\int_0^t \dd W_s\Bigr] = \E[W_t^2] &= t &&\text{(variance grows linearly in time)}, \label{eq:ito:var}\\
  (\dd W_t)^2 &= \dd t &&\text{(Itô's rule; quadratic variation)}. \label{eq:ito:qvar}
\end{align}
\eqref{eq:ito:var} is the continuous-time origin of the $\sqrt T$ scaling in
\eqref{eq:clt}: volatility over a horizon $h$ is $\sigma\sqrt h$, and annualising daily
volatility by $\sqrt{252}$ is the same statement in discrete time.

\subsection{The OU process: definition, solution, and moments}
\label{sec:ou:definition}

\begin{definition}[Ornstein--Uhlenbeck]
\begin{equation}
  \dd X_t \;=\; \theta\bigl(\mu - X_t\bigr)\dd t + \sigma\dd W_t,
  \qquad \theta>0,
  \label{eq:ou}
\end{equation}
with $\theta$ the \emph{speed of mean reversion}, $\mu$ the long-run mean, and $\sigma$
the instantaneous diffusion.
\end{definition}

\paragraph{Solving it.} Multiply by the integrating factor $e^{\theta t}$:
\[
  \dd\bigl(e^{\theta t}X_t\bigr) = e^{\theta t}\theta\mu\dd t + e^{\theta t}\sigma\dd W_t ,
\]
(the $e^{\theta t}(-\theta X_t)\dd t$ term cancels the drift's $X$-dependence), integrate
from $0$ to $t$, divide by $e^{\theta t}$:
\begin{equation}
  X_t \;=\; \mu + \bigl(X_0-\mu\bigr)e^{-\theta t} + \sigma\int_0^t e^{-\theta(t-s)}\dd W_s .
  \label{eq:ou:solution}
\end{equation}
The stochastic integral is Gaussian with mean zero and, by the Itô isometry
$\E[(\int f\dd W)^2]=\int f^2\dd s$,
\begin{equation}
  \Var(X_t) \;=\; \sigma^2\int_0^t e^{-2\theta(t-s)}\dd s
  \;=\; \frac{\sigma^2}{2\theta}\bigl(1-e^{-2\theta t}\bigr).
  \label{eq:ou:variance}
\end{equation}
So the conditional moments are
\begin{equation}
  \E[X_t\mid X_0] = \mu + (X_0-\mu)e^{-\theta t},
  \qquad
  \Var(X_t\mid X_0) = \frac{\sigma^2}{2\theta}\bigl(1-e^{-2\theta t}\bigr),
  \label{eq:ou:moments}
\end{equation}
and as $t\to\infty$ the process converges to its stationary (invariant) distribution
\begin{equation}
  X_\infty \sim \mathcal N\Bigl(\mu,\ \frac{\sigma^2}{2\theta}\Bigr),
  \qquad
  \Cov(X_t, X_{t+\ell}) = \frac{\sigma^2}{2\theta}e^{-\theta\lvert\ell\rvert},
  \label{eq:ou:stationary}
\end{equation}
i.e.\ an AR(1)-style exponentially decaying autocovariance with $\phi=e^{-\theta}$
(compare \eqref{eq:ar1}).

\paragraph{Half-life.} From \eqref{eq:ou:moments}, the expected deviation from $\mu$
halves when $e^{-\theta t}=\tfrac12$:
\begin{equation}
  \boxed{\ \tau_{1/2} \;=\; \frac{\ln 2}{\theta}\ }
  \label{eq:halflife}
\end{equation}
This single number governs the whole strategy's economics: it says how long the
mispricing lives, and hence how long you must hold, how much you can expect to earn per
trade, and how much turnover (and cost) you will incur.

\begin{intuition}
An OU process is a spring with friction, hit by random kicks. $\theta$ is the spring
stiffness: stiff springs snap back fast (small $\tau_{1/2}$, many trades, high cost,
each trade earning little); loose springs drift back slowly (few trades, each earning
more, but you must survive a long exposure). $\sigma$ is the size of the kicks and sets
the stationary dispersion $\sigma/\sqrt{2\theta}$: the typical size of the deviation you
get to trade. The ratio that matters for profit-per-trade is roughly
$\text{(deviation captured)} \times \text{(probability of reversion)} - \text{(cost)}$,
and both terms are functions of $\theta$ and $\sigma$.
\end{intuition}

\subsection{Estimating $\theta,\mu,\sigma$ from daily data}
\label{sec:ou:estimation}

\subsubsection{The discretised regression (what \code{statkit.ou} does)}

The Euler discretisation of \eqref{eq:ou} with sampling interval $\Delta t=1$ is
\[
  X_{t+1}-X_t \;=\; \theta\mu\Delta t - \theta X_t\Delta t + \sigma\sqrt{\Delta t}\,\eta_{t+1},
  \qquad \eta\sim\mathcal N(0,1)\ \text{i.i.d.}
\]
which is an ordinary linear regression of the first difference on a constant and the
lagged level:
\begin{equation}
  \Delta X_t \;=\; a + b\,X_{t-1} + e_t,
  \qquad
  \hat\theta = -\hat b,\quad
  \hat\mu = -\hat a/\hat b,\quad
  \hat\sigma = \hat s_e\ \ (\Delta t=1),
  \label{eq:ou:regression}
\end{equation}
and $\hat\tau_{1/2}=\ln 2/\hat\theta$. The exact discretisation
\eqref{eq:ou:moments} gives the same regression with $b=e^{-\theta}-1\approx-\theta$ for
small $\theta$, so the Euler version is accurate for daily data where
$\theta\approx0.07$; the exact maximum-likelihood version (which regresses $X_t$ on
$X_{t-1}$ and converts $b=e^{-\theta}$) is preferred when $\theta$ is large.

\begin{remark}[Why this is not the same as an ADF test --- and why we run both]
\eqref{eq:ou:regression} with $H_0:b=0$ is algebraically the Dickey--Fuller regression
\eqref{eq:df}. The \emph{estimator} is the same; the \emph{inference} is not. Under
$H_0:b=0$ with a stationary \emph{alternative}, the $t$-statistic has the non-standard
DF distribution \eqref{eq:dfdistribution}, not a normal one --- so you cannot report
``$\hat\theta$ is significant at 5\%'' from a normal table. We therefore report the ADF
$p$-value (correct critical values, via \code{statsmodels}) for the \emph{test}, and the
OU regression for the \emph{magnitude} $\hat\theta$ and $\hat\tau_{1/2}$. Conflating the
two is a common error.
\end{remark}

\subsubsection{Bias, and how good our estimate is}

The OLS estimate of $b$ in an AR(1) is biased downward (toward more mean reversion) in
finite samples by approximately $-(1+3\phi)/T$; with $\phi\approx1$ and $T=4{,}070$ this
is a bias of $\approx-10^{-3}$ in $b$, i.e.\ $\approx-10^{-3}$ in $\theta$. For the
static hedge ($\hat\theta\approx0.004$) that bias is 25\% of the estimate --- the
half-life of $\sim171$ bars is not precise. For the rolling spread
($\hat\theta\approx0.074$) the bias is $\approx1.4\%$, and the half-life of $9.4$ bars is
trustworthy to roughly a bar. This is why Phase~\ref{ch:phase6} reports the half-life
for both periods and both models rather than a single headline number.

The toolkit's own recovery test (\code{scripts/run\_statkit\_tests.py},
Table~\ref{tab:statkit_validation}) simulates an OU process with a \emph{known}
$\theta=0.4$ and estimates $\hat\theta=0.3965$ --- a $0.9\%$ error --- with an implied
half-life of $1.748$ against the true $\ln 2/0.4 = 1.733$ ($0.9\%$). Knowing the truth is
the only way to validate an estimator, which is why the synthetic tests exist.

\subsection{From half-life to a trading rule}
\label{sec:ou:trading}

\subsubsection{How much of a reversion do you capture?}

If you enter at deviation $z_0$ (in standard-deviation units) and hold for $h$ bars, the
expected remaining deviation is $z_0e^{-\theta h}$, so the expected captured reversion is
\begin{equation}
  \E[\text{gain}] \;=\; z_0\bigl(1-e^{-\theta h}\bigr) .
  \label{eq:ou:gain}
\end{equation}
At $h=\tau_{1/2}$ you capture 50\%; at $h=2\tau_{1/2}$, 75\%; at $h=3\tau_{1/2}$,
87.5\%. The marginal gain per extra bar is $z_0\theta e^{-\theta h}$, which decays
geometrically --- there is a natural stopping point, and holding much longer than
$2$--$3$ half-lives buys almost nothing while paying carry, borrow, and tail risk.

\paragraph{The measured match.} Phase~\ref{ch:phase6} finds
$\hat\tau_{1/2}=9.36$--$9.44$ bars for the traded rolling spread and a realized average
holding period of $16.2$ bars (full sample) / $17.3$ bars (2022+): \textbf{1.7--1.9
half-lives per trade}, capturing $\approx70$--$73\%$ of the expected reversion by
\eqref{eq:ou:gain}. That is a reasonable horizon by the theory --- but it sits on the
turnover-heavy side, and turnover is what the 10\,bps cost model punishes. The exit band
$z_{\text{exit}}=0.5$ (rather than 0) is what keeps the average hold at 16 rather than 30
bars; the sensitivity surface of Phase~\ref{ch:phase7} shows that tightening the entry
band and lengthening the window (fewer, longer, better-separated trades) is the single
biggest net-Sharpe lever available, from $0.05$ at $(W{=}120, z_{\text{entry}}{=}2.0)$ to
$0.21$ at $(W{=}180, z_{\text{entry}}{=}3.0)$.

\subsubsection{The optimal-threshold intuition}

A formal treatment: with entry threshold $z_e$, exit $z_x$, per-bar carry $\kappa$, and
round-trip cost $c$, the expected profit per trade is approximately
\begin{equation}
  \pi(z_e) \;\approx\; \sigma_S\Bigl[(z_e-z_x)\bigl(1-e^{-\theta h}\bigr)\Bigr]
  \;-\; c\;-\;\kappa h(z_e),
  \qquad
  h(z_e)\approx \frac{1}{\theta}\log\frac{z_e-z_x}{\text{target}},
  \label{eq:ou:profit}
\end{equation}
and the number of trades per year falls as $z_e$ rises (roughly like the tail probability
of the stationary distribution, $2\Phi(-z_e)$ per bar of exposure). Raising $z_e$
therefore increases profit per trade and decreases trade count. The optimum trades off
the two, and it moves \emph{up} as $c$ rises. This is exactly the pattern in the
measured heatmaps: at 10\,bps the best $(W,z_{\text{entry}})$ cell is
$(180, 3.0)$ with net Sharpe $0.206$, whereas the low-threshold cell $(60,1.5)$ is
$-0.244$. Cost, not signal, is what sets the optimal trading frequency.

\subsection{Why the OU model can be wrong about your spread}
\label{sec:ou:caveats}

Four failure modes, each observed or guarded against in this project:
\begin{enumerate}
  \item \textbf{The mean $\mu$ is not constant.} A rolling hedge makes $\hat\alpha_t$ and
    $\hat\beta_t$ move, so the ``equilibrium'' the spread reverts to drifts. This is the
    KPSS rejection on the full-sample rolling spread (Phase~\ref{ch:phase6}: ADF rejects
    a unit root but KPSS rejects strict level-stationarity, verdict \emph{ambiguous}).
    Remedy: standardise on a trailing window, not the full sample --- which is what
    equation~\eqref{eq:saa:z} does.
  \item \textbf{$\theta$ is state-dependent.} Reversion is faster in calm markets and
    slower (or negative) in crises, when correlations jump and relative-value positions
    diverge. The regime table of Phase~\ref{ch:phase7} shows exactly that: net Sharpe
    $-0.18$ in the lowest volatility tercile, $-0.78$ in the middle, $-0.34$ in the
    highest, and $-1.31$ in the 2022 selloff window.
  \item \textbf{$\sigma$ is not constant.} Volatility clustering
    (Section~\ref{sec:ts:arch}) means the stationary variance $\sigma^2/2\theta$ is
    time-varying, so a fixed $z$-threshold means different things in 2011 and 2020.
    Remedy: the trailing $\hat\sigma_{S,t}$; better (untested here): volatility-targeted
    sizing.
  \item \textbf{The process is not Gaussian.} Real spreads have fat tails and can jump
    (earnings, M\&A, a peer's guidance cut). An OU fit will treat a permanent break as a
    large temporary deviation and keep adding to the position. This is the classic
    ``picking up pennies in front of a steamroller'' failure of pairs trading, and no
    amount of statistical testing on historical data removes it --- only position limits,
    stop-losses, and fundamental monitoring do. Phase~\ref{ch:phase7} implements the
    limits; the monitoring is future work.
\end{enumerate}
